
% ============================================================================
\clearpage
\section*{v4.8.0-r05 Source-Only: Corrected Explicit Formula in the Field-Equation and \texorpdfstring{$\Theta$}{Theta}--\texorpdfstring{$R$}{R} Layer}
\addcontentsline{toc}{section}{v4.8.0-r05 Source-Only: Corrected Explicit Formula in the Field-Equation and Theta--R Layer}
\label{sec:v480-r05-corrected-explicit-formula}
% ============================================================================

\paragraph{Purpose.}
Layer v4.8.0-r04 recorded that the explicit formula \(\Sigma=1-\Theta+R\) taken from
\cite{brendecke2026rh} is misstated (finding A1) and listed the sites that use it.  This layer
carries the corrected formula into those sites: the Einstein-equation reading of
Theorem~\ref{thm:einstein}, the bare-\(\Lambda\) proposition, the \(\Theta\)--\(R\) cancellation
statement, the photon-frame \(\Theta\)-sum conjecture of OP~5 Step~(B), and the GUE theorem.
Two statements can now be decided numerically; both are decided here.  Earlier layers are not
rewritten.

\noindent\textbf{Source-only status.}
This v4.8.0-r05 layer is source-only.  It does not build a PDF, does not assign a new DOI,
and does not authorize a public release.  The v4.8.0 Version DOI
\href{https://doi.org/10.5281/zenodo.23365244}{10.5281/zenodo.23365244} and the Concept DOI
\href{https://doi.org/10.5281/zenodo.19210728}{10.5281/zenodo.19210728} are unchanged by it.

\subsection*{v4.8.0-r05 corrected trace decomposition}
\label{subsec:v480-r05-trace-decomposition}

Write \(\Sigma(\tau):=\sum_\rho e^{i\gamma_\rho\tau}\), summed over all non-trivial zeros
\(\rho=\tfrac12+i\gamma_\rho\) (the zeros are used as numerically known data; no statement about
their real parts is assumed).

\paragraph{Theoremlet v4.8.0-r05.1.}
For \(\tau>0\), in the sense of distributions,
\[
\Sigma(\tau)
=\underbrace{e^{\tau/2}}_{M(\tau)}
\;-\;\underbrace{e^{-\tau/2}\sum_{n\ge2}\Lambda(n)\,\delta(\tau-\log n)}_{P(\tau)}
\;-\;\underbrace{\frac{e^{-\tau/2}}{e^{2\tau}-1}}_{Z(\tau)} .
\]
\(M\) comes from the pole of \(\zeta\) at \(s=1\), \(P\) from the prime powers with von Mangoldt
weights, \(Z\) from the trivial zeros.

\paragraph{Proof sketch.}
This is the classical explicit formula
\(\sum_n\Lambda(n)\delta(\tau-\log n)=e^\tau-\sum_\rho e^{\rho\tau}-\sum_{k\ge1}e^{-2k\tau}\),
multiplied by \(e^{-\tau/2}\), with \(\sum_k e^{-2k\tau}=1/(e^{2\tau}-1)\).  Paired with a smooth
test function supported in \((0.6,3.2)\) and 300 zeros, both sides of the unmultiplied identity
agree to six digits (\(4.393268\)), while the form \(1-\Theta+R\) of the source gives
\(2.1587\) against \(0.5198\) (r04, A1).  \(\square\)

\paragraph{What changes.}
The three-part structure survives: pole, primes, trivial zeros.  What changes is every
functional form.  The constant \(1\) becomes the growing term \(e^{\tau/2}\); the primes enter with
weight \(\Lambda(n)=\log p\) and damping \(e^{-\tau/2}\), not with weight \(1/m\); the trivial-zero
term has denominator \(e^{2\tau}-1\) and positive sign inside \(Z\).  In addition,
equation~\eqref{eq:Theta-def} uses \(\Theta\) for the zero sum itself, which is also the
definition of \(\Sigma\); read literally, \eqref{eq:gw} then gives \(2\Sigma=1+R\).  The symbol
\(\Theta\) is reserved here for the prime-side distribution only.

\subsection*{v4.8.0-r05 the Einstein-equation reading}
\label{subsec:v480-r05-einstein-reading}

Theorem~\ref{thm:einstein} varies \(S=\int_0^\infty\Sigma(\tau)\,d\tau\).  With the corrected
decomposition this integral does not exist: \(M\) grows like \(e^{\tau/2}\), \(Z(\tau)\sim1/(2\tau)\)
as \(\tau\to0^+\), and \(\Sigma\) is a distribution, not a function.  Moreover
\(\mathrm{Tr}\,e^{-i\tau L_0}\) is the trace of a unitary group, not a heat kernel, so the heat-kernel
variation formula used in the proof does not apply.

\paragraph{Theoremlet v4.8.0-r05.2.}
The correspondence
\[
M\ \leftrightarrow\ \Lambda g_{\mu\nu},\qquad
P\ \leftrightarrow\ T_{\mu\nu},\qquad
Z\ \leftrightarrow\ G_{\mu\nu}
\]
is retained as a structural conjecture.  Theorem~\ref{thm:einstein} is not derived: a well-defined
action requires a test-function pairing \(\langle\Sigma,\varphi\rangle\) and an explicit dependence
of \(\varphi\) on \(g^{\mu\nu}\), neither of which is supplied.

\paragraph{Proof sketch.}
Divergence of the three pieces is read off from Theoremlet r05.1; the trace statement follows
from the unitarity of \(e^{-i\tau L_0}\).  \(\square\)

\paragraph{Bare \(\Lambda\).}
Proposition~\ref{prop:bare-lambda} assigns to the constant term a size \(\gamma_1^2\).  In the corrected
formula the pole term is \(M(\tau)=e^{\tau/2}\); it carries no dependence on any zero, so no
\(\gamma_1^2\) can arise from it.  The proposition is withdrawn as a derivation; the scale
\(\gamma_1^2\) remains an unsupported ansatz.

\paragraph{Cancellation.}
Theorem~\ref{thm:partial-cancel} reads the explicit formula as a cancellation of the pole by the
other terms.  The formula is an identity between distributions; it does not by itself fix the value
of a truncated integral.  The open content of Step~(B) is exactly such a value, decided below.

\subsection*{v4.8.0-r05 decision of the \texorpdfstring{$\Theta$}{Theta}-sum conjecture}
\label{subsec:v480-r05-theta-sum-decision}

Proposition~\ref{conj:theta-sum-B} states
\[
\mathcal T(T)=2\sum_{n\ge1}\frac{\sin(T\gamma_n/\gL)}{\gamma_n/\gL}
=\frac{3T}{2}-\log T-\frac{T^2}{24}-\gamma_E-\bsym^2T^2\approx0.8934,
\qquad T=\gL^{-3/2}\approx0.3928 .
\]

\paragraph{Theoremlet v4.8.0-r05.3.}
\(\mathcal T(T)=-0.1097\pm0.0009\).  The conjecture is false; the two sides differ by \(1.003\).
Satisfying the equation would require \(\bsym^2\approx6.80\), outside \([0,1)\).

\paragraph{Proof sketch.}
\(\mathcal T(T)=\gL\,S(a)\) with \(S(a)=2\sum_{\gamma>0}\sin(a\gamma)/\gamma\) and
\(a=T/\gL\approx0.2107<\log2\).  For \(x=e^a\in(1,2)\) one has \(\psi_0(x)=0\), so the
Riemann--von Mangoldt formula gives exactly
\(\sum_\rho x^\rho/\rho=x-\log2\pi-\tfrac12\log(1-x^{-2})\).  Pairing \(\rho\) with \(\bar\rho\),
\[
S(a)=e^{-a/2}\Bigl(x-\log2\pi-\tfrac12\log(1-x^{-2})\Bigr)
-\sum_{\gamma>0}\frac{\cos a\gamma}{\tfrac14+\gamma^2}
+\sum_{\gamma>0}\frac{\sin a\gamma}{2\gamma(\tfrac14+\gamma^2)},
\]
with both remaining sums absolutely convergent.  With the first 2000 zeros the value is
\(-0.109652\) (50 zeros: \(-0.109464\); 500: \(-0.109661\)), and the tail is bounded by
\(8.2\cdot10^{-4}\) through the zero density \((2\pi)^{-1}\log(\gamma/2\pi)\).  The direct partial sums
converge to the same value (\(-0.1449\), \(-0.1117\), \(-0.1053\) at \(N=50,500,2000\)); the
first agrees with the \(-0.146\) recorded in Remark~\ref{rem:theta-sum-structure}.  \(\square\)

\paragraph{Consequence.}
By Remark~\ref{rem:explicit-formula-remainder} the conjecture was stated as equivalent to the
balance condition~\eqref{eq:tr-balance} and to OP~5 Step~(B).  That equivalence used the old
\(R\); with the corrected decomposition it would have to be re-derived.  Either way, the balance
statement in the form given there does not hold.  The defect of almost exactly \(1\) suggests that
the constant term was mishandled in the regularization; even a shift by \(1\) leaves a gap of
\(0.003\), larger than the error bound.  Prediction P4 and the photon-frame \(\Lambda\) programme,
already marked open in r03, lose this route.

\subsection*{v4.8.0-r05 GUE statistics: empirical, not derived}
\label{subsec:v480-r05-gue}

Theorem~\ref{thm:gue} derives GUE statistics from self-adjointness, irreducibility and the
absence of time reversal.  The derivation does not hold: self-adjointness and irreducibility do not
imply any random-matrix class (the harmonic oscillator has neither GUE nor GOE spacings), and an
antiunitary \(J\) with \(J^2=1\) commuting with \(L_0\), as asserted in the proof, would be a
time-reversal-type symmetry and would point to GOE, not GUE.

\paragraph{Theoremlet v4.8.0-r05.4.}
The nearest-neighbour spacings of the zeros follow GUE statistics empirically.  For zeros
\(101\)--\(2000\), unfolded with \(N(t)=\tfrac{t}{2\pi}\log\tfrac{t}{2\pi e}+\tfrac78\), the
Kolmogorov--Smirnov distance to the Wigner surmise is \(0.041\) for GUE (\(\beta=2\)), \(0.108\)
for GOE (\(\beta=1\)) and \(0.315\) for Poisson; small spacings \(s<0.1\) occur with frequency
\(0.0005\) (GUE surmise \(0.0011\), GOE \(0.0078\), Poisson \(0.095\)).

\paragraph{Proof sketch.}
Direct computation from 2000 zeros.  The residual GUE distance reflects the Wigner surmise as an
approximation and the low height; it is consistent with the Montgomery--Odlyzko results.
\(\square\)

Theorem~\ref{thm:gue} is therefore retained as an empirical statement about the zeros and
withdrawn as a consequence of the Sphaera structure.

\subsection*{v4.8.0-r05 status of the affected statements}
\label{subsec:v480-r05-status}

\begin{center}
\small
\begin{tabular}{@{}>{\raggedright\arraybackslash}p{0.34\textwidth}>{\raggedright\arraybackslash}p{0.58\textwidth}@{}}
\toprule
statement & status after r05\\
\midrule
\eqref{eq:gw}, \eqref{eq:Theta-def}, \eqref{eq:R-def} & replaced by Theoremlet r05.1\\
Thm.~\ref{thm:einstein} & structural conjecture (r05.2); action not defined\\
Prop.~\ref{prop:bare-lambda} & withdrawn as derivation; \(\gamma_1^2\) unsupported\\
Thm.~\ref{thm:partial-cancel} & identity, not cancellation; content moved to Step~(B)\\
Prop.~\ref{conj:theta-sum-B}, Rem.~\ref{rem:explicit-formula-remainder} & false: \(\mathcal T(T)=-0.1097\), not \(0.893\)\\
Thm.~\ref{thm:gue} & empirical fact (r05.4); derivation withdrawn\\
\bottomrule
\end{tabular}
\end{center}

\subsection*{v4.8.0-r05 ledger and guard update}
\label{subsec:v480-r05-ledger-guard}

\[
R_{\rm fe}^{v480\text{-}r05}=(0,0,1,1,0)
\]
for (corrected trace decomposition; \(\Theta\)-sum conjecture decided; Einstein-equation reading
not derived; photon-frame balance condition without a valid form; GUE statement settled
empirically).

The layer spends no guard:
\[
R_{\rm glob}^{(v480\text{-}r05,+)}=(0,1,1,1,1),
\qquad
\Gamma_{\rm guard}^{v480\text{-}r05}=(1,1,1,1,1),
\qquad
\beta_4=1/256\ \text{unspent}.
\]
It does not imply exact finite \(C5\) zero closure, exact all-readout terminality,
\(C6/T2\), a global photon-clock/RH theorem, or global all-readout closure.

